% Primitive: bayes-rule — authored from probability/bayes_rule_manim.py and
% its README concepts table (six scenes), plus plan 004's verified anchors.
% Backward-grounded: the conditional chapter supplies the front-door
% identity, the cohort counts and the CI license; forward pointers live in
% "Where this goes" only. Problem answers are computed by
% answers/bayes_rule.py (the solve-gate anchor); new anchor keys ride
% primitives/anchors-proposals-probability.yaml until merged.
\chapter{Bayes' rule}

\begin{narrative}
The conditional chapter ended one visible step short of something: the
shared-numerator identity $P(A)\,P(B \mid A) = P(B)\,P(A \mid B)$,
standing on the page with nothing yet asked of it. This chapter asks.
One division turns the identity into the law of updating beliefs — and
then the real work begins, because the formula is the least useful form
of the law. Counts compute it, odds compress it, and a game-show host
stress-tests what it actually takes as input.

\section{Through the front door}

Divide the standing identity by $P(B)$ — positivity inherited from the
conditional chapter's definitions — and Bayes' rule is already on the
page:
\[
  P(A \mid B) \;=\; \frac{P(B \mid A)\,P(A)}{P(B)} .
\]
Nothing new was proved; the division is the two-slices picture re-read,
and the denominator is the conditional chapter's law of total
probability run over the hypothesis columns. What is new is the
\emph{reading}, and the three names that come with it: $P(A)$ is the
\emph{prior} (what the hypothesis was worth before the news),
$P(B \mid A)$ is the \emph{likelihood} (how well the hypothesis
explains the news), $P(A \mid B)$ is the \emph{posterior} (what it is
worth after). The claim of the whole chapter, in one line: evidence
does not \emph{determine} beliefs, it \emph{updates} them —
posterior $\propto$ prior $\times$ likelihood, normalize last.

\section{Counting it out}

The first real computation needs no formula at all. Diseasitis, in
whole students: a $100$-student cohort, $20$ sick and $80$ healthy.
The test flags $\anchor{004.diseasitis.sickpos}$ of the sick and
$\anchor{004.diseasitis.healthypos}$ of the healthy. A student tests
positive; the probability they are sick is the share of positives who
are sick:
\[
  \frac{18}{18 + 24} \;=\; \anchor{004.diseasitis.posterior} .
\]
That is Bayes' rule, computed by counting — the prior never appeared
because it travelled \emph{inside} the counts ($20$ sick students out
of $100$ is the prior, already multiplied in). This is why natural
frequencies cure base-rate neglect where the probability format
fails~\cite{probability-gigerenzer-hoffrage-1995-how-to-improve-bayesian-reasoning-w}:
posed in probabilities, the problem is solved by
$\anchor{004.natfreq.lay}$ of laypeople — and reformatting alone moves
physicians from $\anchor{004.natfreq.gyn}$. When you must
\emph{communicate} a posterior, hand over the counts.

\section{The odds form}

Now compress. Keep only the \emph{ratio} between the two hypotheses
and watch what the evidence does to it — the waterfall:

\begin{center}
\begin{tikzpicture}[scale=0.9]
  % Scene 3's device: stream widths are the prior odds, the surviving
  % fraction of each stream is the likelihood, the pool is the
  % posterior odds. Sick Cool, healthy Muted, blocked water Warm,
  % the pool (the result) outlined Accent.
  \fill[Cool!40] (0,4.4) rectangle (1,5.2);
  \fill[Muted!30] (1.5,4.4) rectangle (5.5,5.2);
  \node[above, color=CoolInk] at (0.5,5.55) {sick $1$};
  \node[above, color=Muted] at (3.5,5.55) {healthy $4$};
  \fill[Cool!40] (0,2.6) rectangle (0.9,4.4);
  \fill[Warm!30] (0.9,4.1) rectangle (1,4.4);
  \fill[Muted!30] (1.5,2.6) rectangle (2.7,4.4);
  \fill[Warm!30] (2.7,4.1) rectangle (5.5,4.4);
  \node[right, color=Muted] at (3.15,3.3) {pass-through $3{:}1$};
  \draw[->, Muted] (0.45,2.35) -- (1.75,1.85);
  \draw[->, Muted] (2.1,2.35) -- (2.85,1.85);
  \fill[Cool!40] (1.6,0.6) rectangle (2.5,1.6);
  \fill[Muted!30] (2.5,0.6) rectangle (3.7,1.6);
  \draw[AccentInk, thick] (1.6,0.6) rectangle (3.7,1.6);
  \node[right, color=AccentInk] at (4.1,1.1)
    {$\anchor{004.diseasitis.postodds} \;\to\;
      \anchor{004.diseasitis.posterior}$};
\end{tikzpicture}
\end{center}

Two streams pour down at the prior widths, $1$ sick to $4$ healthy.
The test lets $9$ in $10$ of the sick stream through and $3$ in $10$
of the healthy stream — pass-through in ratio $3{:}1$ — and the pool
at the bottom collects
$\anchor{004.diseasitis.sickpos}{:}\anchor{004.diseasitis.healthypos}$,
which is $\anchor{004.diseasitis.postodds}$: the same
$\anchor{004.diseasitis.posterior}$ as the counted cohort. The
waterfall is the cohort chips, and the tree, and the two slices, drawn
as one object with renormalization deferred to the very end — and its
law is a single
multiplication~\cite{probability-yudkowsky-so8res-smith-et-al-bayes-rule-odds-form}:
\[
  \text{posterior odds} \;=\; \text{prior odds} \times
  \text{likelihood ratio} .
\]
Only ratios matter: scale either stream's width or both pass-through
fractions by any common factor and the pool's ratio never moves. One
discipline comes with the compression — odds and probabilities answer
different questions, so keep $\anchor{004.diseasitis.postodds}$ and
$\anchor{004.diseasitis.posterior}$ visibly distinct.

\section{One test, two patients}

The odds form factors the conditional chapter's prevalence pair into
its parts. That test — $9$ in $10$ sensitive, $1$ in $10$ false
positives — has \emph{one} number of its own: the likelihood ratio
$\anchor{004.LR}$~\cite{probability-3blue1brown-the-medical-test-paradox}.
The posterior belongs to the patient:
\[
  1{:}9 \times 9 = 1{:}1 \;\to\; \tfrac12,
  \qquad
  1{:}99 \times 9 = \anchor{004.lowprev.odds} \;\to\;
  \anchor{003.prevalence.low} .
\]
The counted $\anchor{003.prevalence.highcounts}$ and
$\anchor{003.prevalence.lowcounts}$ of the conditional chapter fall
out as derivations instead of tallies — same answers, one
multiplication each. And the marketing phrase ``$90\%$ accurate''
stands revealed as one word hiding two numbers (sensitivity and false
positive rate), being silently read as a third (the posterior). A
posterior can never be stated without its prior; a test's quality can
never be stated as one number that \emph{is} a posterior.

\section{Yesterday's posterior}

Evidence arrives in sequences, and the odds form makes iteration
effortless: yesterday's posterior is today's prior. On the conditional
chapter's own trick coins — heads $9$ in $10$ versus $1$ in $10$, so
each head carries likelihood ratio $\anchor{004.LR}$ — start at
$1{:}1$: one head moves the odds to $9{:}1$, a second to
$\anchor{004.coins.HH}$; a head \emph{then a tail} lands back at
exactly $1{:}1$~\cite{probability-mit-18-05-class-11-bayesian-updating}.
That exact cancellation is the point: if evidence \emph{replaced}
belief, the order and mix of observations could not wash out; because
evidence \emph{reweights} belief, a tail's ratio $\tfrac19$ undoes a
head's $9$ automatically. Two captions guard the habit. Likelihood
ratios multiply only when the observations are conditionally
independent \emph{given the hypothesis} — the conditional chapter's
closing lesson, cashed at the exact moment it is needed (it is the
disease being common cause of both test results that makes the second
test still informative, yet lets its ratio simply multiply on). And a
prior of zero stays zero under any evidence of nonzero probability —
no finite likelihood ratio rescues a hypothesis that was never in the
running.

\section{The host's protocol}

Monty Hall, at last — deferred by the conditional chapter until the
tools could handle it honestly. You pick door 1; the host, who knows
where the car is, opens door 3 on a goat; switch to door 2? The trap
is conditioning on the revealed \emph{fact} (``door 3 has a goat'')
instead of the host's \emph{action} under his protocol — the
two-children lesson at series scale. Put a uniform prior on the car's
position and let the likelihood be the probability of the action
\emph{the host actually took}, one row per
protocol~\cite{probability-rosenthal-monty-hall-monty-fall-monty-crawl}:
the standard host (never the car, coin-flip when free) assigns the
opening likelihoods $(\tfrac12, 1, 0)$ across the three car positions,
and switching wins with probability $\anchor{004.monty.standard}$;
Monty Fall, who slipped and opened a random door that merely
\emph{happened} to hide a goat, assigns
$(\tfrac12, \tfrac12, 0)$ — and switching wins only
$\anchor{004.monty.fall}$; Monty Crawl, who opens the lowest-numbered
door he can, was \emph{forced} high in opening door 3 — and switching
wins with certainty. Same door, same goat, three answers: the
information is not in what was revealed but in what the protocol made
that revelation \emph{cost}. Rosenthal's proportionality principle —
posterior $\propto$ prior $\times$ likelihood of what was observed —
is just Bayes with the uniform prior left implicit, and it is the
series' closing rule: \emph{condition on what happened, the way it
happened — then multiply}.

\paragraph{Where this goes.} The multiplications of this chapter are
about to become additions: the logarithm chapters build a ruler on
which a likelihood ratio is a fixed \emph{distance}, and iterated
updating is walking — with a zero prior sitting infinitely far down
the scale, which is why no finite evidence reaches it. The likelihood,
named here on the front door, returns as the lens the softmax chapter
reads per-frame distributions through. And at the road's end, the
decoding chapter cashes the independence chapter's common-cause
warning the Bayesian way: an external language model is a prior,
multiplied onto per-frame likelihoods — condition on the way the
audio happened, then multiply.
\end{narrative}

\section*{Practice problems}

\begin{problem}
In the Diseasitis cohort — $100$ students, $20$ sick of whom the test
flags $9$ in $10$, $80$ healthy of whom it falsely flags $3$ in $10$ —
a student tests positive. Compute the probability they are sick using
whole-person counts only.
\begin{hint}
How many sick positives; how many positives in all?
\end{hint}
\begin{solution}
Sick positives: $\tfrac{9}{10} \cdot 20 =
\anchor{004.diseasitis.sickpos}$. Healthy positives:
$\tfrac{3}{10} \cdot 80 = \anchor{004.diseasitis.healthypos}$. The
positives total $42$, and the share of them that are sick is
$18/42 = \anchor{004.diseasitis.posterior}$. No formula was used — the
prior ($20$ sick in $100$) rode inside the counts.
\end{solution}
\end{problem}

\begin{problem}
Recompute the Diseasitis posterior in the odds form, and check it
against the counts.
\begin{hint}
Prior odds from $20{:}80$; the likelihood ratio from the two flag
rates.
\end{hint}
\begin{solution}
Prior odds: $20{:}80 = 1{:}4$. Likelihood ratio:
$\tfrac{9/10}{3/10} = 3$. Posterior odds:
$1{:}4 \times 3 = \anchor{004.diseasitis.postodds}$ — literally the
$18{:}24$ of the counted cohort — and converting odds to probability,
$\tfrac{3}{3+4} = \anchor{004.diseasitis.posterior}$. Two routes, one
answer; the waterfall is the cohort with renormalization deferred.
\end{solution}
\end{problem}

\begin{problem}
A test has sensitivity $\tfrac{9}{10}$ and false positive rate
$\tfrac{1}{10}$. Using the odds form, find $P(\text{sick} \mid +)$ for
a patient from a population at $10\%$ prevalence, and for one at
$1\%$.
\begin{hint}
The test contributes one number; only the prior odds change.
\end{hint}
\begin{solution}
The likelihood ratio is $\tfrac{9/10}{1/10} = \anchor{004.LR}$. At
$10\%$ prevalence the prior odds are $1{:}9$, so posterior odds
$1{:}9 \times 9 = 1{:}1$, i.e.\ $\tfrac12$. At $1\%$ the prior odds
are $1{:}99$, so $1{:}99 \times 9 = 9{:}99 =
\anchor{004.lowprev.odds}$, i.e.\ $\anchor{003.prevalence.low}$.
These reproduce the conditional chapter's counted
$\anchor{003.prevalence.highcounts}$ and
$\anchor{003.prevalence.lowcounts}$ exactly — the test supplied one
number, and each patient supplied their own prior.
\end{solution}
\end{problem}

\begin{problem}
One of two coins (heads $\tfrac{9}{10}$ versus $\tfrac{1}{10}$) is
chosen by a fair flip. Track the odds that you hold the heads-heavy
coin after observing H, after HH, and after HT — and give each as a
probability.
\begin{hint}
One likelihood ratio per flip; multiply. What is a tail's ratio?
\end{hint}
\begin{solution}
Each head multiplies the odds by $\tfrac{9/10}{1/10} = \anchor{004.LR}$;
each tail by $\tfrac{1/10}{9/10} = \tfrac19$. From $1{:}1$: after H,
$9{:}1$, probability $\tfrac{9}{10}$; after HH,
$\anchor{004.coins.HH}$, probability $\tfrac{81}{82}$; after HT,
$9 \times \tfrac19 = 1$, exactly $1{:}1$ — probability $\tfrac12$, as
if nothing had been seen. Full conditioning over both coins and all
flip sequences gives the same three posteriors, which is the licensed
multiplication at work: the flips are conditionally independent given
the coin.
\end{solution}
\end{problem}

\begin{problem}[stretch]
You pick door 1; the host opens door 3, showing a goat. Compute the
probability that switching to door 2 wins under each protocol:
(a) the standard host — never opens your door or the car's, coin-flips
when both goat doors are available; (b) Monty Fall — opened a random
door of the other two and happened to reveal a goat; (c) Monty Crawl —
always opens the lowest-numbered door he legally can.
\begin{hint}
Uniform prior over the car's position; the likelihood is the
probability of \emph{this host} opening door 3, given each car
position.
\end{hint}
\begin{solution}
Prior $\tfrac13$ each. The likelihood of ``opens door 3'' per car
position $(1, 2, 3)$: (a) standard — $(\tfrac12, 1, 0)$, so the
posterior on door 2 is
$\tfrac{1}{1 + 1/2} = \anchor{004.monty.standard}$; (b) Fall —
$(\tfrac12, \tfrac12, 0)$, posterior
$\anchor{004.monty.fall}$: the accidental reveal carries less
information because it was equally cheap either way; (c) Crawl — he
opens door 2 whenever it is legal, so opening door 3 forces the car
onto door 2: switching wins with certainty. Same revealed fact, three
answers — the likelihood belongs to the \emph{action under the
protocol}, not to the fact.
\end{solution}
\end{problem}
